\scanpage{6}
\diagram{\begin{tikzcd}[ampersand replacement=\&]
C\times_S B\arrow[r,"{\psi\times \mathrm{id}_B}"]\arrow[d,"\mathrm{pr}_1"']\arrow[dr,phantom,"{\text{cart}}"{description,draw,ellipse,inner sep=1pt}] \& B^*\times_S B^{\deleted{*}}\arrow[d,"\mathrm{pr}'_1"]\\
C\arrow[r,"\psi"'] \& B^*
\end{tikzcd}}
Ceci dit, on aura
\[
\deleted{R \mathrm{pr}_{2*}(N)}\qquad N\simeq(\psi\times \mathrm{id}_B)^*(W_B)
\]
($W_B$ le Module de Weil sur $B\times_S B^*$), d'où
\[
\begin{aligned}
\mathbb{R} \mathrm{pr}_{2*}(N)&\simeq \mathbb{R} \mathrm{pr}_{2*}(\psi\times \mathrm{id}_B)^*(W_B)\\
&\overset{\text{chgt de base}}{\simeq}\mathbb{L}\psi^*(\mathbb{R} \mathrm{pr}'_{2*}(W_B)).
\end{aligned}
\]

\uline{Lemme 1} Soit $B$ un schéma abélien sur $S$\inserted{, de dim relative $n$}, $W$ le faisceau de Weil sur $B^*\times_S B^{\deleted{\unsure{*}}}$, $p\colon B^*\times_S B^{\deleted{\unsure{*}}}\unsure{\mathrel{-}|\to}B^*$ la projection. Alors
\[
\mathbb{R} p^{\deleted{\unsure{*}}}_*(W)\simeq e_{\deleted{\unsure{B}}*}\bigl[\deleted{\unsure{\mathcal{O}\otimes\tau[-n]_{B^*}}}\tau_{B^*}[-n]\bigr)
\]
(isom. canonique), où $\tau_{B^*}=\bigwedge^n\underline{t}_{B^*}$ ($t_{B^*}$ le faisceau tangent à $B^*$ le long de la section $0$).
\margin{$\simeq e_*(\omega^{-1})$ où $\omega=\det^{*}(\underline{\Omega}^1_{B/\unsure{A|S}})$}

On en conclut un

\uline{Corollaire} Soient $C,B$ schémas abéliens sur $S$\inserted{, $B$ de dim relative $n$}, $N$ un Module inversible birigidifié sur $C\times_S B$, $\psi\colon C\longrightarrow B^*$ l'hom correspondant. Alors, on a
\[
\mathbb{R} \mathrm{pr}_{1*}(N)\simeq \mathbb{L}\psi^*(\unsure{\tau|e|e_*\tau}_{B^*}\deleted{\unsure{*}}[-n])\quad\deleted{\unsure{a}}
\]
En particulier, si $\psi$ est un épimorphisme, \deleted{\unsure{et si|donc}} donc plat, et si $K=\operatorname{Ker}\psi$, alors $K$ est un sous-schéma en groupe de $C$ plat de prés. finie sur $S$, donc \deleted{\unsure{définit un|défini} \unsure{complexe}} $\underline{O}_K$ est parfait sur $C$, et on a
\[
\mathbb{R} \mathrm{pr}_{1*}(N)\simeq\underline{O}_K[-n]\otimes_S\tau_{B^*}
\]

Dans le cas qui nous occupe, $C=A\times A$, \deleted{$\unsure{C^*\simeq A}$} $\varphi\colon A\simeq C^*$, et \deleted{\unsure{est}} par suite $K$ est l'\uline{antidiagonale} de $A\times_S A$, noyau de $\pi_A\colon A\times_S A\longrightarrow A$, qui est le graphe de la \uline{symétrie} $s_A$ de $A$. Donc \deleted{\unsure{la} \illegible{1}} \unsure{on gagne}, modulo dém. du Lemme 1.
% archive sheet number: 5
